\chapter{Conclusions}
\label{chapter:conclusions}

\section{Summary}

This report describes the first phase of the development of an instrumented
canine prosthesis for objective locomotion monitoring.
The motivation is grounded in a clear clinical gap: prosthetic adaptation in
animals is currently assessed through subjective observation, which limits
consistency and reproducibility.

The proposed approach addresses this gap by embedding inertial and force
sensing directly into the prosthetic structure, enabling continuous,
non-invasive, and quantitative assessment during real-world activity.
A review of the state of the art confirmed that no existing solution combines
kinetic and kinematic sensing in a single-limb wearable prosthetic device.

\section{Achievements}

The main achievements of this phase are:

\begin{itemize}
  \item A fully functional hardware prototype integrating the MPU-9250 IMU,
        HX711 load cell, microSD storage, and ESP32-C3 wireless
        microcontroller in a compact, battery-powered form factor;
  \item A firmware pipeline operating at 80\,Hz, including force filtering,
        quaternion-based orientation estimation, local CSV storage, and a
        browser-accessible web interface for real-time monitoring and data
        retrieval;
  \item Sensor validation confirming correct peripheral initialisation, load
        cell calibration, and orientation estimation with smooth dynamic
        tracking;
  \item A preliminary dataset of 11 acquisition runs recorded with the
        prototype fitted to a canine subject in an initial field session,
        demonstrating that the system captures the periodic structure of
        the gait cycle in both force and inertial channels under real conditions.
\end{itemize}

These results demonstrate that the proposed instrumented prosthesis is viable
as a platform for biomechanical data acquisition, and provide a solid
foundation for the planned work outlined in Chapter~\ref{chapter:workplan}.

\section{Next Steps}

The immediate priorities for the next phase are: structured data acquisition
sessions with canine subjects under controlled and free-movement conditions;
IMU mounting calibration to align the sensor frame with the prosthetic
anatomical frame; and feature extraction from the collected dataset to begin
characterising adaptation-related patterns.
